\chapter*{Preface: Two Verbs and a Choice}
\addcontentsline{toc}{chapter}{Preface: Two Verbs and a Choice}

This book leans on two verbs you will not find in a dictionary. It uses them on nearly every page, and it is fairer
to set them down here, at the door, than to let the reader guess at them for a hundred pages. The verbs are
\emph{tange} and \emph{funge}.

To \emph{tange} is to select by a characteristic: to reach into a collection and pick out \emph{this one, not that
one}, on the strength of some property that tells the two apart. You tange a ripe apple from a bin --- not any apple,
that one, because of what it is. To \emph{funge} is the opposite gesture: to bag together by a like characteristic, to
gather into one heap all the things that are, in the way that matters, the same. You funge your coins by denomination,
dropping every quarter into the same roll, because for the purpose at hand one quarter is as good as another. Tange
divides by a difference; funge unites by a sameness. Between the two of them they are the whole of what it means to
sort a world.

They are also, this book will argue, the whole of what it means to \emph{measure}. To measure a thing is first to
tange it --- to select it out of everything it is not, by the one characteristic the instrument can read --- and then
to funge it --- to declare it the same as everything the instrument cannot tell it apart from. A ruler tanges a length
out of the continuum and funges together every length that falls between its two nearest marks. The marks are the
funge; the reading is the tange; and the little gap between two marks is exactly the set of lengths the ruler has
quietly decided are, for its purposes, one and the same. Every measurement in these pages is built from those two acts
and nothing else, and the reader who keeps both verbs in hand will find the whole construction easier to walk for it.

\medskip

There is a place where \emph{tange} stops being an ordinary word and becomes one of the deepest questions in
mathematics, and it is worth walking up to that place now --- because the book walks up to it and then, on purpose,
declines to cross. Tange one thing from one bag and no one raises an eyebrow. Tange one thing from each of finitely
many bags --- a coin from each of ten rolls --- and still no one objects; you can simply do it, one bag after the
next, and be finished in ten steps. But now ask to tange one thing from each of \emph{infinitely} many bags, all at
once, with no rule that says which one to take from each --- and you have asked for something that cannot be taken for
granted. That such a selection can always be made, even across infinitely many bags, even with no rule to make it by,
is a separate assumption with a name of its own: the \emph{axiom of choice}. The axiom of choice is \emph{tange},
promoted to the infinite and granted for free.

That axiom has a short and stormy history, and it earns a paragraph here because this book's whole discipline is a
reply to it. Zermelo~\cite{zermelo1904} stated it plainly in 1904, in order to prove that every set can be
well-ordered --- lined up so that every part of it has a first element. The proof went through, but the tool alarmed
people, because it asserted that a choice could always be made while never once saying \emph{how}. Borel, Lebesgue,
and Baire --- mathematicians who wanted every object exhibited and not merely posited --- objected that a choice with
no rule behind it was not a choice at all but a wish. The alarm grew louder when Banach and Tarski~\cite{banachtarski1924} showed, in 1924, that the axiom lets you cut a solid ball into a handful of pieces and reassemble them,
by rigid motions alone, into \emph{two} balls each the size of the first --- a paradox that is no sleight of hand but a
theorem, and a standing warning about how much strange power an unrestricted choice quietly carries. For three decades
the deeper question hung open: is the axiom even consistent with the rest of mathematics, and is it truly needed?
G\"odel~\cite{godel1938} settled the first half in 1938 --- the axiom can introduce no contradiction that was not
already present without it. Cohen~\cite{cohen1963} settled the second in 1963, with a method he called \emph{forcing}:
he built a mathematical universe in which all the other axioms hold and the axiom of choice fails, proving choice
\emph{independent} --- neither compelled nor forbidden by the rest. And Cohen's forcing did its work by approaching an
ungraspable object through an ascending tower of \emph{finite} conditions, each one a thing you could actually write
down, converging on a generic limit that no single condition ever reached. The reader will meet that shape again
before long; it is the shape of this entire construction.

\medskip

For this book takes a side, and the side is the finite one. The machine at the center of these pages tanges and funges
without pause --- it is doing nothing else --- but it only ever tanges \emph{finitely}, and only ever from bags it can
genuinely tell apart. It never reaches into infinitely many bags at once; it never selects by a rule it cannot state;
it never posits a choice it cannot exhibit. When it names a thing, it does so by counting the thing into one of
finitely many boxes, and a choice among finitely many boxes needs no axiom at all --- you can just look. The
consequence is exact, and it is checkable: the construction is \emph{choice-free}. Ask the machine to audit what its
results actually lean on, and the axiom of choice is not among them; the only assumptions that survive the audit are
the plainest ones ordinary logic already grants. The book does not avoid the axiom of choice as a matter of taste. It
earns every selection it makes, one finite decision at a time, and it can prove that it did.

That discipline is not an incidental virtue; it is the reason the book can end where it does. An instrument that
helped itself to an infinite, ruleless choice could posit any answer it pleased and dress the wish as a result --- it
could hand you the exact value of the number this book reaches for and dare you to object. This instrument cannot, and
will not. It tanges only what it can distinguish and funges only what it genuinely cannot tell apart; and when it runs
clean out of distinctions, it stops, and hands back not a fabricated point but an honest interval --- a bracket around
the answer, the width of the bracket the exact measure of what its finite choices could not resolve. The two verbs are
the method; the axiom of choice is what the first of them becomes when it is handed the infinite for free; and the
whole of this book's honesty lies in its refusal of that gift. It tanges, it funges, and it stays finite --- and
because it stays finite, the answer it finally gives is a thing it earned rather than a thing it wished for.
